% ============================================================
%  Altar Valley Berms — Figure Registry Report  (Paper 1)
%  Paper 1: Berm condition and vegetation response
%  Auto-generated from figures/paper1/figure_registry_concise.txt
%  Date: 2026-02-28
% ============================================================
\documentclass[11pt, a4paper]{article}

\usepackage[utf8]{inputenc}
\usepackage[T1]{fontenc}
\usepackage{lmodern}
\usepackage[left=2.5cm, right=2.5cm, top=2.5cm, bottom=2.5cm]{geometry}
\usepackage{graphicx}
\usepackage{tabularx}
\usepackage{booktabs}
\usepackage{caption}
\usepackage{float}
\usepackage{xcolor}
\usepackage{parskip}
\usepackage{microtype}
\usepackage{hyperref}
\hypersetup{colorlinks=true, linkcolor=black, urlcolor=blue}
\usepackage{csvsimple}
\usepackage{longtable}

% ── column types ─────────────────────────────────────────────
\newcolumntype{L}{>{\raggedright\arraybackslash}X}   % wrapping left col

% ── helper: shaded row label ─────────────────────────────────
\newcommand{\rowlbl}[1]{\textbf{\small #1}}

% ── figure + registry table macro ────────────────────────────
%   \figentry{label}{file}{updated}{stats text}{interp text}
\newcommand{\figentry}[5]{%
  \subsection*{#1}
  \begin{figure}[H]
    \centering
    \includegraphics[width=\linewidth]{../figures/paper1/#2}
    \caption*{\textit{File:} \texttt{#2} \quad \textit{Updated:} #3}
  \end{figure}
  \vspace{0.5em}
  \begin{tabularx}{\linewidth}{@{} p{3.2cm} L @{}}
    \toprule
    \rowlbl{Field} & \rowlbl{Content} \\
    \midrule
    \rowlbl{Figure} & #1 \\[2pt]
    \rowlbl{File} & \texttt{\small #2} \\[2pt]
    \rowlbl{Updated} & #3 \\[6pt]
    \rowlbl{Statistical test} & #4 \\[6pt]
    \rowlbl{Interpretation} & #5 \\
    \bottomrule
  \end{tabularx}
  \vspace{3em}
}

% ============================================================
\begin{document}

\begin{center}
  {\LARGE \textbf{Geomorphic and Soil Controls on Berm Structural Integrity and Vegetation Response in the Altar Valley, Arizona}} \\[0.3em]
  {\large Figure Registry Report} \\[0.3em]
  {\small Generated 2026-02-28}
\end{center}

\vspace{1em}
\hrule
\vspace{1.5em}

This document summarises each registered figure for Paper 1: the figure image is
shown above a two-column table giving the filename, last-updated date, the
statistical test applied, and a brief interpretation of the key result.

\vspace{1em}

% ── Fig 1 ────────────────────────────────────────────────────
%% FIG_1_START
\figentry%
  {fig1}%
  {fig1_rf_permutation_importance.png}%
  {2026-03-02 15:31}%
  {Random Forest (n_estimators=300, class_weight='balanced'), 5-fold stratified CV. Permutation importance = mean decrease in holdout AUC from shuffling each predictor (30 repeats). Top 5 predictors shown for each outcome. CV AUC (condition) = 0.73 ± 0.01 (n=743). CV AUC (vegetation response) = 0.65 ± 0.04 (n=743).}%
  {Flow accumulation and berm length dominate condition prediction. Hillslope gradient and soil properties (texture, typical profile) are more important for vegetation response than for condition, reflecting different mechanisms.}
%% FIG_1_END

% ── Fig 1 (alternate) ────────────────────────────────────────
%% FIG_1_ALT_START
\figentry%
  {fig1_alt}%
  {fig1_alt.png}%
  {2026-03-02 15:32}%
  {Drop-column (leave-one-out) importance: Random Forest (n_estimators=300, class_weight='balanced'), 5-fold stratified CV. Importance = decrease in mean CV AUC when predictor is removed from the full model. Top 5 predictors shown per outcome. Full-model CV AUC: condition = 0.73 ± 0.01, vegetation response = 0.65 ± 0.04.}%
  {Drop-column importances corroborate permutation importance rankings. Removing flow accumulation or berm length produces the largest AUC drop for condition; soil properties and landform show the largest effect for vegetation response.}
%% FIG_1_ALT_END

% ── Fig 2 ────────────────────────────────────────────────────
%% FIG_2_START
\figentry%
  {Figure 2}%
  {fig2\_geomorph\_soil\_intact\_effective.png}%
  {2026-02-27 08:45}%
  {Panel A: Pearson $\chi^2$ ($\varphi$ coefficient) for a
   2$\times$2 Intact $\times$ Vegetation-response contingency table ($N = 743$).
   Panel B: two-sided Fisher's exact test comparing proportions between Intact and
   Degraded berms within each soil texture class.}%
  {Berm intactness and vegetation response are statistically independent
   ($\varphi = -0.01$, $p = 0.739$), indicating that intact and degraded berms
   produce vegetation responses at similar overall rates. The exception is Sandy
   loam, where intact berms show significantly higher vegetation response than
   degraded berms ($p = 0.034$).}
%% FIG_2_END

% ── Fig 3 ────────────────────────────────────────────────────
%% FIG_3_START
\figentry%
  {Figure 3}%
  {fig3\_berm\_condition.png}%
  {2026-02-27 09:39}%
  {Two-category panels (length, slope) use a one-sided Fisher's exact test;
   multi-category panels (landform, texture) use pairwise Fisher's exact tests
   with the overall $\chi^2$ $p$-value shown in the panel title.}%
  {Berm condition is strongly associated with berm length --- shorter berms are
   markedly more likely to be intact ($p < 0.001$) --- and modestly associated
   with slope class ($p < 0.05$). Landform does not significantly predict
   condition ($p = 0.204$), but soil texture does ($p < 0.001$).}
%% FIG_3_END

% ── Fig 4 ────────────────────────────────────────────────────
%% FIG_4_START
\figentry%
  {fig4}%
  {fig4_vegetation_response.png}%
  {2026-03-02 15:29}%
  {Two-category panels (length, slope) use a one-sided Fisher's exact test; multi-category panels (landform, texture) use pairwise Fisher's exact tests with the overall chi-squared p-value shown in the panel title.}%
  {Vegetation response is significantly associated with slope class (p < 0.001), landform (p < 0.001), and soil texture (p < 0.001), but not with berm length alone. Steep-slope berms and flood-plain berms show the highest rates of vegetation response.}
%% FIG_4_END

% ── SI Fig 1 ─────────────────────────────────────────────────
%% SI_FIG_1_START
\figentry%
  {Supplementary Figure 1}%
  {SI\_fig1\_fa30\_pointplot\_by\_landform.png}%
  {2026-03-02 15:29}%
  {Point plots show mean $\pm$ 95\% CI of 30\,m flow accumulation (log scale) by landform, grouped by vegetation response (Panel A) and structural condition (Panel B).}%
  {Flood-plain berms experience the highest contributing flow accumulation regardless of outcome, while fan-terrace berms receive the least. Within each landform, effective berms tend to have similar or higher flow accumulation than ineffective berms, and degraded berms tend toward higher values than intact berms, consistent with greater hydraulic stress on degraded structures.}
%% SI_FIG_1_END

% ── SI Fig 2 ─────────────────────────────────────────────────
%% SI_FIG_2_START
\figentry%
  {Supplementary Figure 2}%
  {SI\_landform\_by\_predictors.png}%
  {2026-02-27 08:43}%
  {Two-sided binomial test ($H_0$: 50/50 split) applied within each landform
   group for each predictor category (length, slope, clay content, soil
   development); labels: *** $p < 0.001$, ** $p < 0.01$, * $p < 0.05$,
   ns~$p \geq 0.05$.}%
  {Predictors are non-randomly distributed across landforms --- high-clay soils
   concentrate on flood plains, steep slopes on fan terraces, and B-horizon
   development on fan and stream terraces --- confirming that landform is a
   meaningful stratification variable and motivating its inclusion as a covariate.}
%% SI_FIG_2_END

\newpage
\section*{Supplementary Tables}

\subsection*{Supplementary Table 1: Univariate GLM Rankings for Vegetation Response Predictors}

These tables present univariate predictor rankings using binomial generalized linear models (GLMs) to identify the most informative individual predictors for each outcome. For each predictor, we report:
\begin{itemize}
  \item \textbf{McFadden pseudo-R\textsuperscript{2}}: measures improvement in log-likelihood relative to the null model (higher values indicate better fit)
  \item \textbf{Tjur R\textsuperscript{2}}: measures separation in predicted probabilities between outcome classes
  \item \textbf{CV AUC}: 5-fold cross-validated area under the ROC curve (higher values indicate better predictive performance)
  \item \textbf{LRT p-value}: likelihood ratio test comparing the model to the null model
\end{itemize}

These univariate rankings complement the multivariate Random Forest analysis (Figure 1) by showing how each predictor performs in isolation, helping to identify the strongest individual associations before accounting for interactions and non-linear effects.

\vspace{1em}

\begin{table}[H]
\centering
\caption*{\textbf{Top 15 predictors for Vegetation Response}}
\small
\csvreader[
  tabular=lrrrr,
  table head=\toprule \textbf{Predictor} & \textbf{McFadden $R^2$} & \textbf{Tjur $R^2$} & \textbf{CV AUC} & \textbf{LRT $p$} \\ \midrule,
  late after line=\\,
  table foot=\bottomrule
]{../figures/paper1/SI_table_predictors_vegetation_response_clean.csv}
{predictor=\pred,mcfadden_r2=\mcfad,tjur_r2=\tjur,cv_auc=\cvauc,lrt_p=\lrtp}
{\pred & \mcfad & \tjur & \cvauc & \lrtp}
\end{table}

\vspace{2em}

\subsection*{Supplementary Table 2: Univariate GLM Rankings for Structural Integrity Predictors}

\begin{table}[H]
\centering
\caption*{\textbf{Top 15 predictors for Structural Integrity}}
\small
\csvreader[
  tabular=lrrrr,
  table head=\toprule \textbf{Predictor} & \textbf{McFadden $R^2$} & \textbf{Tjur $R^2$} & \textbf{CV AUC} & \textbf{LRT $p$} \\ \midrule,
  late after line=\\,
  table foot=\bottomrule
]{../figures/paper1/SI_table_predictors_structural_integrity_clean.csv}
{predictor=\pred,mcfadden_r2=\mcfad,tjur_r2=\tjur,cv_auc=\cvauc,lrt_p=\lrtp}
{\pred & \mcfad & \tjur & \cvauc & \lrtp}
\end{table}

%% SI_FIG_3_START
\figentry%
  {Supplementary Figure 3}%
  {SI\_fig3\_countplots\_by\_outcome\_and\_landform.png}%
  {2026-03-02 15:29}%
  {3$\times$3 panel of count plots. Rows: soil development (B horizon vs.\ none), berm length class (short vs.\ long), slope class (shallow vs.\ steep). Columns: vegetation response, structural condition, landform.}%
  {Fan-terrace berms are overwhelmingly found in B-horizon soils, whereas stream-terrace and flood-plain berms dominate in soils without a B horizon. Berm length and slope class are more evenly distributed across outcomes, though short berms are disproportionately intact and steep-slope berms show a higher proportion of vegetation response.}
%% SI_FIG_3_END

\end{document}
